\chapter{Два появления \texorpdfstring{$1/7$}{1/7}: стабильный \texorpdfstring{$K_0$}{K0}-мост и предел каноничности}

Предыдущий гейт обнаружил нормированную суперразмерность
$(20-15)/35=1/7$. Ранее тот же вес возник у единственной голономной
семейной линии как $(15/35)(1/3)=1/7$. Настоящая глава проверяет, является
ли это численным совпадением, одной и той же следовой бухгалтерией или
операторным мостом между двумя секторами.

Сначала необходимо исправить возможное слишком быстрое чтение. Семейный
проектор $P_0$ имеет ранг один в $\mathbb C^3_{\rm fam}$; он не является
проектором ранга пять внутри $M_{35}$. Поэтому прямое равенство
$[P_0]=[p_{20}]-[p_{15}]$ в $K_0(M_{35})$ не типизировано. Сравнение
требует общей стабилизированной алгебры.

\section{Две точные формулы}

Связывающая градуировка даёт
\begin{equation}
 \tau_{35}(p_{20}-p_{15})
 =\frac{20-15}{35}=\frac17.
 \label{eq:v5-one-seventh-corner}
\end{equation}
Голономный нулевой проектор
\begin{equation}
 P_0=\frac13(I+C_3+C_3^2),
 \qquad \rank P_0=1,
\end{equation}
в наблюдаемом углу имеет вес
\begin{equation}
 \tau_{35}(p_{15})\tau_3(P_0)
 =\frac{15}{35}\frac13=\frac17.
 \label{eq:v5-one-seventh-holonomy}
\end{equation}
Равенство дробей эквивалентно целочисленному тождеству
\begin{equation}
 (20-15)\cdot3=15\cdot1=15.
 \label{eq:v5-one-seventh-rank-duality}
\end{equation}
Следовательно, совпадает не десятичное приближение, а произведение двух
независимо замороженных разложений: разности углов $20-15$ и семейного
разложения $3=1+2$.

\section{Общая сравнительная алгебра}

Для корректной типизации введём только как бухгалтерский контейнер
\begin{equation}
 \mathcal B_{\rm cmp}=M_{35}(\mathbb C)\otimes M_3(\mathbb C)
 \simeq M_{105}(\mathbb C).
 \label{eq:v5-one-seventh-comparison-algebra}
\end{equation}
Тогда виртуальный угловой класс стабилизируется как
\begin{equation}
 ([p_{20}]-[p_{15}])\otimes[I_3]
\end{equation}
и имеет виртуальный ранг $5\cdot3=15$. Голономный класс
\begin{equation}
 [p_{15}\otimes P_0]
\end{equation}
имеет положительный ранг $15\cdot1=15$. Поэтому в
$K_0(\mathcal B_{\rm cmp})\simeq\mathbb Z$ выполняется
\begin{equation}
 \boxed{
 ([p_{20}]-[p_{15}])\otimes[I_3]
 =[p_{15}\otimes P_0]
 }
 \label{eq:v5-one-seventh-k0-identity}
\end{equation}
при стандартной идентификации классов по рангу.

\begin{proposition}[Стабильная природа повторения $1/7$]
Два появления $1/7$ являются одним точным равенством нормированных следов
и одним равенством стабильных $K_0$-рангов после помещения в
$M_{105}$. Поэтому их нельзя считать независимыми численными
предсказаниями, но также нельзя считать случайным совпадением несвязанных
вещественных чисел.
\end{proposition}

Важно, что $\mathcal B_{\rm cmp}$ пока не объявляется новой физической
родительской алгеброй. Она лишь помещает оба уже существующих чтения в
одну типизированную группу.

\section{Почему равенство классов ещё не является переносом}

Левая часть \eqref{eq:v5-one-seventh-k0-identity} виртуальна. Чтобы
представить её положительным проектором до тензорирования, требуется
выбрать ранг-пять подпространство
\begin{equation}
 r_5\le p_{20}.
\end{equation}
Но известные центральные редукции, совместимые с замороженными
$J$-парами и целыми семейными блоками, имеют ранги
\begin{equation}
 \{0,2,6,8,12,14,18,20\}.
 \label{eq:v5-one-seventh-known-ranks}
\end{equation}
Канонического ранга пять в этом списке нет.

Произвольный выбор существует: пространство всех таких подпространств
есть
\begin{equation}
 \operatorname{Gr}(5,20),
 \qquad
 \dim_{\mathbb R}\operatorname{Gr}(5,20)=2\cdot5\cdot15=150.
\end{equation}
После выбора $r_5$ проекторы $r_5\otimes I_3$ и
$p_{15}\otimes P_0$ имеют одинаковый ранг пятнадцать, поэтому между ними
существует частичная изометрия $V$. Однако при фиксированных концах выбор
$V$ образует торсор по $U(15)$ вещественной размерности $225$.

\begin{nogocriterion}[Равенство $K_0$ не выбирает сплетающий оператор]
Ранг и нормированный след доказывают существование стабильной
эквивалентности, но не выбирают ни положительный представитель $r_5$, ни
частичную изометрию $V$, ни её совместимость с $J$, калибровочным
действием и голономией. Подстановка произвольного элемента
$\operatorname{Gr}(5,20)$ или $U(15)$ снова была бы ручным секторным
выбором.
\end{nogocriterion}

\section{Что именно теперь требуется вывести}

Следующий объект должен быть не новой численной формулой, а граничным
гомоморфизмом или трансгрессией
\begin{equation}
 \partial_{\rm hol}:
 [p_{15}\otimes P_0]
 \longmapsto
 ([p_{20}]-[p_{15}])\otimes[I_3],
 \label{eq:v5-one-seventh-transgression-target}
\end{equation}
возникающей из уже имеющихся хопфовой линии, расширения Мориты или
относительной точной последовательности. Такой механизм должен выполнять
сразу четыре условия:
\begin{enumerate}
 \item не выбирать базис в $\mathbb C^{20}$;
 \item быть ковариантным относительно семейной голономии и $J$;
 \item порождать положительный представитель виртуального класса;
 \item фиксировать частичную изометрию до калибровочной эквивалентности.
\end{enumerate}

Если такой граничный морфизм существует, постоянная пятимерная
необратимость угла может оказаться не удаляемым фоном, а тенью
голономного класса после трансгрессии. Если он не существует, равенство
$1/7$ останется точной, но только следовой двойственностью.

\section{Вердикт}

Повторение $1/7$ впервые получило строгое общее выражение. Оно связывает
разность углов и голономную линию в стабильном $K_0$, однако не даёт
оператора переноса. Поэтому предыдущая отрицательная проверка нечётного
ядра не отменена, но финальная заморозка откладывается ровно на один
топологический тест граничной трансгрессии.

\begin{protocol}
\begin{enumerate}
 \item угловое чтение $1/7$: \textbf{точно};
 \item голономное чтение $1/7$: \textbf{точно};
 \item общая целочисленная формула: \textbf{$5\cdot3=15\cdot1$};
 \item равенство стабильных $K_0$-рангов в $M_{105}$:
 \textbf{получено};
 \item равенство является независимым численным предсказанием:
 \textbf{нет};
 \item канонический положительный представитель ранга пять в $H_{20}$:
 \textbf{не получен};
 \item каноническая частичная изометрия: \textbf{не получена};
 \item операторный мост: \textbf{не построен};
 \item физическое замыкание: \textbf{не достигнуто};
 \item следующий гейт:
 \textbf{граничная трансгрессия класса $1/7$}.
\end{enumerate}
\end{protocol}

Машинный сертификат сохранён в
\path{s2t_v5_one_seventh_k0_bridge_gate_results.json}.